%======================================================================
% فصل چهارم: پیاده‌سازی
%======================================================================
\chapter{پیاده‌سازی}
\label{ch:implementation}

\section{مقدمه}

در این فصل، نحوه‌ی تبدیل طرح مفهومی فصل قبل به یک سامانه‌ی اجرایی توضیح
داده می‌شود. پیاده‌سازی انجام‌شده صرفاً یک نمونه‌ی آزمایشی نیست؛ بلکه
یک کدپایه‌ی عملیاتی است که لایه‌های مدیریتی، مالی و درگاهی را در کنار
هم ارائه می‌کند. در این پیاده‌سازی، تلاش شده است که هم از امکانات آماده‌ی
چارچوب \lr{Laravel} استفاده شود و هم منطق اختصاصی پروژه در لایه‌های مستقل
و آزمون‌پذیر قرار گیرد.

\section{فناوری‌های استفاده‌شده}

\subsection{چارچوب \lr{Laravel}}

هسته‌ی سمت سرور سامانه با چارچوب \lr{Laravel} پیاده‌سازی شده است
~\cite{laravel2024}. انتخاب این چارچوب به این دلیل انجام شد که هم از
معماری لایه‌ای و الگوهای متداول توسعه به‌خوبی پشتیبانی می‌کند، هم ابزارهای
استانداردی برای اعتبارسنجی، صف، کش، رویداد و احراز هویت در اختیار
می‌گذارد. افزون بر آن، پشتیبانی داخلی از \lr{Pipeline} با نیاز اصلی پروژه
هم‌راستا است و سازگاری آن با \lr{SQL Server} و \lr{Redis} در محیط
ویندوز، پیاده‌سازی عملیاتی سامانه را ساده‌تر کرده است.

افزون بر این، استفاده از \lr{Laravel Sanctum} برای نقاط پایانی مدیریتی،
احراز هویت نشست‌محور کاربر را ساده کرده است~\cite{sanctum2024}.

\subsection{پایگاه داده‌ی \lr{SQL Server}}

ذخیره‌سازی داده‌های پایدار با \lr{SQL Server} انجام شده است~\cite{sqlserver2024}.
این پایگاه داده محل نگهداری اطلاعات کاربران، کلیدها، ارائه‌دهندگان،
مدل‌ها، کیف پول، اشتراک‌ها، فاکتورها، گزارش‌ها و لاگ‌های ممیزی است.
استفاده از تراکنش‌های پایگاه داده برای عملیات مالی، مانند شارژ و برداشت
از کیف پول، یکی از الزامات اجرایی این بخش بوده است.

\subsection{\lr{Redis}}

\lr{Redis} در این پروژه سه نقش اصلی دارد~\cite{redis2024}. از یک سو برای
کش داده‌های پرتکرار مانند طرح‌ها، ارائه‌دهندگان و پروفایل کاربر به‌کار
می‌رود، از سوی دیگر شمارنده‌های محدودسازی نرخ برای \lr{OTP} و درگاه هوش
مصنوعی را نگه می‌دارد، و در نهایت بستر اجرای صف کارهای پس‌زمینه‌ای مانند
ارسال \lr{OTP} و ارسال گزارش خارجی را فراهم می‌کند.

ترکیب این سه نقش در یک فناوری واحد، پیچیدگی عملیاتی سامانه را کاهش داده
و هماهنگی میان بخش‌های مختلف را بیشتر کرده است.

\subsection{\lr{Vue.js}}

برای پیاده‌سازی \term{میدان بازی}{Playground} از \lr{Vue.js} و ابزار
\lr{Vite} استفاده شده است~\cite{vuejs2024}. این رابط کاربری وظیفه‌ی اصلی
سامانه را بر عهده ندارد، اما برای آزمون دستی نقاط پایانی و مشاهده‌ی سریع
رفتار درگاه بسیار مفید است. رابط مذکور امکان انتخاب نقطه‌ی پایانی،
ارائه‌دهنده، مدل، کلید \lr{API} و محتوای \lr{JSON} را برای کاربر فراهم می‌کند.

\section{ساختار پوشه‌بندی پروژه}

کدپایه به‌صورت حوزه‌محور و لایه‌محور سازمان‌دهی شده است. مهم‌ترین
شاخه‌های پروژه به‌گونه‌ای انتخاب شده‌اند که هر حوزه جای مشخص خود را داشته
باشد. برای مثال، \lr{app/Domains} منطق دامنه‌های اصلی مانند درگاه، احراز
هویت و ارائه‌دهندگان را در خود جای داده است؛ \lr{app/Services} سرویس‌های
سطح کاربرد برای مالی، گزارش، کاربران، کش و مدیریت ارائه‌دهندگان را
نگه می‌دارد؛ \lr{app/Repositories} لایه‌ی دسترسی به داده را از منطق تجاری
جدا می‌کند؛ \lr{app/Http/Controllers} و \lr{app/Http/Resources} به‌ترتیب
وظیفه‌ی کنترل ورودی و شکل‌دهی پاسخ‌های نسخه‌بندی‌شده‌ی \lr{API} را بر عهده
دارند؛ \lr{app/Pipelines} محل خط لوله‌های درگاه و احراز هویت است؛
\lr{app/Jobs} کارهای صف‌شونده‌ای مانند ارسال \lr{OTP}، تولید \lr{PDF}
فاکتور و ارسال گزارش خارجی را نگه می‌دارد؛ و در نهایت
\lr{resources/js/playground} به کد رابط کاربری میدان بازی اختصاص یافته است.

این سازمان‌دهی باعث می‌شود که کدهای مرتبط با یک حوزه‌ی مسئله در کنار هم
قرار گیرند و در عین حال، مسئولیت‌های اجرایی مختلف از هم جدا بمانند.

\section{پیاده‌سازی خط لوله‌ی سامانه‌ی واسط}

هسته‌ی اجرایی سامانه در سرویس \code{GatewayService} قرار دارد. این سرویس
یک شیء \term{زمینه‌ی درخواست}{Gateway Request Context} ایجاد می‌کند و آن را
از چهارده مرحله‌ی خط لوله عبور می‌دهد. این مراحل از نظر منطقی به چند بخش
قابل فهم تقسیم می‌شوند. در ابتدای مسیر، \code{ResolveApiKeyPipe} و
\code{EnforceIpAllowlistPipe} هویت و مجاز بودن مبدأ را بررسی می‌کنند.
پس از آن، مرحله‌های اعتبارسنجی داده و محدودسازی نرخ وارد عمل می‌شوند.
در گام بعد، \code{SelectProviderPipe} و
\code{ValidateProviderSelectionPipe} ارائه‌دهنده و مدل را بارگذاری و صحت
آن‌ها را کنترل می‌کنند. سپس سامانه در دو مرحله‌ی مجزا مصرف تقریبی را
برآورد کرده و وضعیت مالی درخواست را پیش از فراخوانی بیرونی می‌سنجد.

پس از تماس با ارائه‌دهنده، پاسخ در یک گام یکدست‌سازی می‌شود و در انتهای
مسیر، مراحل اندازه‌گیری مصرف، ثبت لاگ، اعمال هزینه و ارسال غیرهمگام گزارش
ساخت‌یافته به‌ترتیب اجرا می‌شوند.

یکی از مزایای عملی این ساختار آن است که در صورت بروز خطا، هر مرحله می‌تواند
زنجیره را متوقف و پاسخ استاندارد تولید کند. به‌این‌ترتیب، خطاهای امنیتی،
مالی یا اعتبارسنجی قبل از تماس با ارائه‌دهنده تشخیص داده می‌شوند.

\section{پیاده‌سازی آداپتورهای ارائه‌دهندگان}

\subsection{\lr{OpenAI}}

آداپتور \lr{OpenAI}-سازگار برای سه ارائه‌دهنده‌ی \lr{OpenAI}، \lr{Groq} و
\lr{OpenRouter} استفاده شده است. این آداپتور با نگاشت داخلی نقاط پایانی
\lr{responses}، \lr{chat/completions}، \lr{embeddings}، \lr{images/generations}،
\lr{audio/transcriptions} و \lr{audio/speech} به مسیرهای متناظر، درخواست را
به ارائه‌دهنده‌ی بالادستی ارسال می‌کند.

در این بخش چند نکته‌ی اجرایی اهمیت دارد. برای
\lr{audio/transcriptions} از ارسال چندبخشی و الحاق فایل استفاده می‌شود،
برای \lr{audio/speech} امکان دریافت پاسخ دودویی و بازگرداندن مستقیم آن
وجود دارد، و مصرف توکن و در صورت امکان تعداد تصاویر از پاسخ استخراج
می‌شود تا در صورت‌حساب‌گذاری به‌کار رود.

\subsection{\lr{Gemini}}

پشتیبانی از \lr{Gemini} در این پروژه از نوع نگاشت بومی است، نه صرفاً تغییر
\lr{Base URL}. آداپتور \lr{GeminiProviderAdapter} درخواست‌های
\lr{chat.completions} و \lr{responses} را به قالب \lr{generateContent} و
درخواست‌های \lr{embeddings} را به قالب \lr{embedContent} یا
\lr{batchEmbedContents} تبدیل می‌کند.

این آداپتور در عمل چند وظیفه‌ی مهم را انجام می‌دهد: ساختار
\lr{messages} یا \lr{input} را به \lr{contents/parts} تبدیل می‌کند،
پیام \lr{system} را جدا کرده و به \lr{systemInstruction} تبدیل می‌نماید،
فراساختار مصرف \lr{Gemini} را به فیلد \lr{usage} در قالب سازگار با
\lr{OpenAI} نگاشت می‌کند و خطاهای بومی \lr{Gemini} را به پاسخ‌های
استاندارد درگاه برمی‌گرداند.

در نسخه‌ی فعلی، \lr{Gemini} برای سه قابلیت گفتگو، \lr{responses} و
جاسازی متن پشتیبانی می‌شود و سایر نقاط پایانی چندرسانه‌ای از مسیر
\lr{OpenAI}-سازگار عبور می‌کنند.

\subsection{\lr{Groq} و \lr{OpenRouter}}

\lr{Groq} و \lr{OpenRouter} با همان آداپتور سازگار با \lr{OpenAI}
پشتیبانی می‌شوند. تفاوت اصلی در این‌جا در سطح پیکربندی نهفته است:
\lr{ProviderRegistry} با توجه به نام ارائه‌دهنده، \lr{Base URL} و کلید
محرمانه‌ی مناسب را از پیکربندی یا پایگاه داده بارگذاری می‌کند. این طراحی
نشان می‌دهد که برای ارائه‌دهندگان سازگار، توسعه‌ی سامانه بیشتر مبتنی بر
داده است تا کدنویسی مجدد.

\section{پیاده‌سازی احراز هویت \lr{OTP}}

منطق احراز هویت کاربر در سرویس \code{OtpService} متمرکز شده است. مرحله‌ی
\term{شروع چالش}{Start} با ایجاد یک \code{OtpContext} و عبور آن از یک
خط لوله‌ی سه‌مرحله‌ای اجرا می‌شود؛ \code{ThrottleOtpPipe} محدودیت نرخ را
بر اساس کانال، مقصد و \lr{IP} اعمال می‌کند، \code{GenerateOtpPipe} کد
عددی تصادفی را تولید و هش کرده و چالش را ثبت می‌کند، و
\code{DispatchOtpPipe} ارسال کد را از طریق \code{SendOtpJob} به صف
می‌سپارد.

در مرحله‌ی \term{تأیید}{Verify}، جدیدترین چالش مربوط به مقصد بازیابی شده و
انقضا، تعداد تلاش‌ها و صحت کد بررسی می‌شود. در صورت موفقیت، کاربر موجود
بازیابی یا در صورت نیاز به‌طور خودکار ایجاد می‌شود. پیاده‌سازی فعلی
\code{SendOtpJob} نقش یک مؤلفه‌ی موقت را دارد و عملیات ارسال واقعی را به ثبت در
لاگ محدود می‌کند؛ بنابراین منطق احراز هویت کامل است، اما اتصال به
ارائه‌دهنده‌ی واقعی پیامک یا ایمیل به‌عنوان کار آتی باقی مانده است.

\section{پیاده‌سازی مدیریت کلیدهای \lr{API}}

مدیریت کلیدهای \lr{API} در دو لایه پیاده‌سازی شده است. در لایه‌ی دامنه،
سرویس \code{ApiKeyService} کلید خام را با پیشوند \code{gk\_} تولید می‌کند،
فقط یک‌بار آن را به کاربر نمایش می‌دهد و سپس نسخه‌ی هش‌شده‌ی آن را
ذخیره می‌کند. علاوه بر هش کامل، یک پیشوند کوتاه نیز نگهداری می‌شود تا
جست‌وجوی کلید در زمان احراز هویت سریع‌تر انجام شود.

در لایه‌ی کاربرد، عملیات ایجاد، ابطال و چرخش کلید با ثبت لاگ ممیزی همراه
است. در سمت دریافت درخواست نیز میان‌افزار \code{AuthenticateApiKey} کلید
را از سرآیند \lr{Authorization: Bearer} یا \lr{X-API-Key} استخراج کرده و
وجود و اعتبار کلید، ابطال یا انقضای آن، فعال بودن مشتری و کاربر مرتبط و
مجاز بودن \lr{IP} را در صورت تعریف فهرست سفید کنترل می‌کند.

این طراحی از ذخیره‌ی رمزهای خام در پایگاه داده جلوگیری کرده و سطح امنیتی
مناسبی برای درگاه فراهم می‌کند.

\section{پیاده‌سازی سیستم صورت‌حساب}

سیستم مالی سامانه از چند بخش مکمل تشکیل شده است. نخست، سرویس
\code{UsageEstimationService} مصرف احتمالی درخواست را برآورد می‌کند. در این
سرویس، متن درخواست‌های گفتگو و \lr{responses} را با استفاده از
\term{توکن‌ساز \lr{BPE}}{Byte Pair Encoding (BPE)} شمارش می‌کند، در صورت
ناشناخته بودن مدل یک \lr{Encoding} پیش‌فرض به‌کار می‌گیرد، برای تصویر
تعداد تصاویر را از فیلدهایی مانند \code{n} استخراج می‌کند و برای صوت
مدت زمان یا تعداد نویسه‌ها را به‌عنوان معیار قیمت‌گذاری غیرتوکنی در نظر
می‌گیرد.

سپس \code{UsageMeteringService} این مقادیر را با \code{pricing\_config}
مدل انتخاب‌شده ترکیب کرده و رکوردهای هزینه‌ای مانند \code{tokens\_in}،
\code{tokens\_out}، \code{images}، \code{audio\_seconds} و
\code{audio\_chars} را می‌سازد.

برای برداشت از کیف پول، \code{WalletService} از تراکنش پایگاه داده
استفاده می‌کند تا افزایش و کاهش موجودی به‌صورت اتمی انجام شود. سرویس
\code{SubscriptionService} نیز منطق ایجاد و لغو اشتراک را به‌همراه برداشت
هزینه‌ی طرح از کیف پول پیاده‌سازی می‌کند. در سطح فاکتور، سرویس
\code{InvoiceService} اقلام را آماده کرده، مجموع و مالیات را محاسبه می‌کند
و فایل \lr{PDF} را از طریق \code{Dompdf} تولید می‌نماید. تولید \lr{PDF}
قابل صف‌شدن است و توسط \code{GenerateInvoicePdfJob} انجام می‌شود.

\section{پیاده‌سازی گزارش‌دهی و ارسال گزارش خارجی}

گزارش‌گیری داخلی و ارسال گزارش خارجی دو مسیر مکمل اما مستقل دارند. در
مسیر داخلی، \code{PersistLogsPipe} در صورت فعال بودن پیکربندی، یک رکورد در
جدول \lr{gateway\_requests} و رکوردهای مصرف متناظر در جدول
\lr{usage\_records} ثبت می‌کند. سپس \code{ReportingService} با تکیه بر این
داده‌ها، گزارش مصرف، دفترکل کیف پول و گزارش فاکتورها را تولید می‌نماید.

برای گزارش‌های تحلیلی بلندمدت، کار \code{AggregateDailyUsageJob} مصرف هر
روز را بر اساس کاربر، کلید، ارائه‌دهنده، مدل و نوع معیار تجمیع می‌کند.
این کار هزینه‌ی گزارش‌گیری در بازه‌های زمانی طولانی را کاهش می‌دهد.

در مسیر خارجی، \code{GatewayLogService} از زمینه‌ی درخواست یک بارِ گزارش
ساخت‌یافته تولید می‌کند، فیلدهای حساس مانند کلید و توکن را حذف یا
\term{پوشانده}{Redact} می‌کند، در صورت نیاز اندازه‌ی داده را می‌کاهد و سپس
از طریق \code{DispatchGatewayLogJob} آن را به یکی از مقاصد پشتیبانی‌شده
می‌فرستد. انتخاب مقصد توسط \code{LogSinkManager} و تنها با تغییر یک مقدار
پیکربندی انجام می‌شود. در نسخه‌ی فعلی دو مقصد \lr{Grafana Loki}
~\cite{grafanaloki2024} و \lr{Better Stack}~\cite{betterstack2024}
پیاده‌سازی شده‌اند.

\section{داشبورد \lr{Playground}}

رابط \lr{Playground} در مسیر \lr{/playground} در اختیار کاربر قرار دارد و
برای آزمون تعاملی نقاط پایانی طراحی شده است. کاربر در این رابط می‌تواند
یکی از شش قابلیت اصلی درگاه را انتخاب کند، نام ارائه‌دهنده و مدل و کلید
\lr{API} را وارد کند، \lr{JSON} درخواست را مستقیماً ویرایش نماید و پاسخ
نهایی را در قالب متنی ببیند.

اگرچه این رابط جایگزین یک پنل مدیریتی کامل نیست، اما برای توسعه و
اشکال‌زدایی بسیار مفید است؛ زیرا می‌تواند بدون نیاز به ابزار بیرونی،
رفتار درگاه و سازگاری پاسخ‌ها را نشان دهد.

در کنار این رابط، مستندات \lr{OpenAPI} نیز با استفاده از
\lr{L5-Swagger} برای پروژه تولید شده است~\cite{l5swagger2024}. ترکیب
مستندات ماشینی و محیط آزمون تعاملی، قابلیت استفاده و نگهداری سامانه را
بهبود داده است.
